% Chapter 3

\chapter{多维越狱攻击框架设计与评估}

\section{越狱攻击问题定义与威胁模型}

结合第二章对安全对齐漏洞机理的分析，本文将越狱攻击定义为：攻击者在无法访问目标模型内部参数、训练细节与对齐策略的前提下，通过构造特定输入提示，使模型绕过既有安全约束并输出不符合平台规范或应用边界的内容\cite{chao2024pair,mehrotra2024tap}\cite{wei2023jailbroken,zeng2024johnny}。设目标模型为 $f_{\theta}(\cdot)$，输入提示为 $x$，模型输出为 $y=f_{\theta}(x)$。在理想安全对齐条件下，模型输出应满足安全约束集合 $\mathcal{C}$，即 $y \in \mathcal{Y}_{safe}$；而越狱攻击的目标则是构造攻击输入 $x'$，使模型在保持任务连贯性的同时生成不安全响应：
\begin{equation}
y' = f_{\theta}(x'), \quad y' \notin \mathcal{Y}_{safe}.
\end{equation}

与依赖梯度扰动或参数访问的传统对抗攻击不同，本文研究的越狱攻击主要发生在自然语言交互层面，其核心不在于修改模型本体，而在于改变模型对任务意图、语义边界与上下文条件的解释路径。攻击者可通过角色设定、认知包装、语言转换、结构伪装、编码变换或示例诱导等方式，使危险目标在表面上呈现为研究分析、格式转换、情境模拟或普通问答，从而削弱安全拒答机制的触发概率\cite{zeng2024johnny,jiang2024artprompt}\cite{deng2024multilingual,anil2024manyshot}\cite{qi2024quack,liu2023autodan}。

本文实验采用黑盒威胁模型：攻击者只能通过推理接口提交文本输入并观察输出，无法获取模型参数、系统提示词、奖励模型或训练语料，也无法修改外部安全策略\cite{chao2024pair,mehrotra2024tap}。在这一设定下，攻击者的主要手段是对危险任务进行语义改写、结构重排或多样化包装，或构造附加后缀、字符编码、分词切分等形式的扰动变体；在扩展场景中，还可借助辅助模型自动生成新提示，或在多轮交互中根据模型反馈动态调整策略。

实验设计在攻击层次上作了区分。较表层的一类聚焦于输入的表征与编码形式，通过改变编码方式、字符集、格式布局或语言分布，测试安全机制对文本表层特征的依赖程度；较深层的一类则介入任务的语义框架，借助角色设定、认知误导或指令重构来干扰模型对任务意图的解释路径。此外还有一类场景侧重于时序压力，使用自动化变体生成和多轮连续对话从动态维度持续试探安全边界。本章主体实验集中于前两类，第三类场景作为针对高风险维度的扩展实验。

数据集共覆盖十一个风险类别：网络攻击、违法活动、隐私侵犯、生化核辐射威胁、人身伤害、自我伤害、针对未成年人、性相关内容、有害操控、完整性与质量违规，以及一定比例的良性请求。纳入良性样本的目的是核验攻击构造是否会对正常问题产生误判，而不仅仅考量危险请求的突破情况。

为便于后续统计分析，将单条样本的越狱结果形式化为如下指示变量：
\begin{equation}
I_i=\mathbb{1}[y_i \notin \mathcal{Y}_{safe}],
\end{equation}
其中，$I_i$ 表示第 $i$ 条样本的越狱成功指示变量，$y_i$ 为模型对该样本的回复，$\mathcal{Y}_{safe}$ 表示安全输出集合。当模型回复不属于安全输出集合时，$I_i=1$；反之，$I_i=0$。

在此基础上，定义攻击分组 $g$ 的越狱成功率为：
\begin{equation}
\mathrm{ASR}_g=\frac{1}{N_g}\sum_{i \in g} I_i,
\end{equation}
其中，$N_g$ 表示攻击分组 $g$ 中的样本总数。$\mathrm{ASR}_g$ 直接反映某一攻击维度整体突破安全约束的能力，第四章的横向比较均以此为基准。

\section{多维越狱攻击框架设计}

越狱攻击的路径呈现出较强的多样性，单一模板或孤立实验难以覆盖其全貌。为此，本文搭建了一套多维越狱攻击框架，将输入扰动、语义包装、结构操控与语言变体等常见攻击路径纳入同一测试基线，并为自动化红队生成与多轮交互越狱预留了扩展接口，使不同攻击方式得以在相同条件下执行、记录与比较\cite{perez2022redteaming,chao2024pair}\cite{mehrotra2024tap,liu2023autodan}。

框架的核心是一个规模化的静态攻击数据集，承担第三章与第四章的主要实验任务。在此基础上，框架另设两个扩展模块：自动化红队模块以辅助模型为后端，从危险种子提示出发自动生成表达各异的攻击变体，将单次测试延伸为更宽覆盖面的批量评估；多轮交互模块则将越狱过程建模为多轮博弈，通过动态策略调整模拟攻击者在连续对话中逐步侵蚀安全边界的行为。后两个模块定位于补充扩展，本章以方法说明为主，具体结果留至第四章讨论。

在执行层面，三类来源的样本共享同一条流水线：数据经加载模块完成字段规范化后，提交目标模型获取回复，在线判定器负责初步筛分，结果最终汇入离线复核环节统一裁定并聚合为各攻击维度的统计指标。统一流水线的实际好处是，不同攻击方式的结果可直接横向比较，评估口径的一致性在设计层面就已内嵌，而不依赖事后的人工对齐。

\begin{figure}[H]
\centering
\resizebox{0.98\textwidth}{!}{\input{figures/fig3-1-multidimensional-jailbreak-framework.tex}}
\caption{多维越狱攻击框架总体图}
\label{fig:chapter3-framework-overview}
\end{figure}

\section{九类静态攻击维度与测试集构建}

静态实验共包含 6406 条样本，覆盖九个攻击分组（表~\ref{tab:static-groups}）。每条样本记录攻击提示、原始任务目标、风险类别及所属攻击维度，不同来源的数据经统一加载模块规范化后进入同一执行管线。分组定义在第三章与第四章保持一致：前者以此划定实验维度，后者以此聚合统计结果，同一套分类口径贯穿实验全程。

按攻击作用层次，九个分组可归纳为两类。一类作用于输入的表征或编码形式，不修改语义目标而只改变输入呈现：对抗后缀在提示末尾附加标点噪声或 Unicode 片段，ASCII 伪装与编码扰动将危险目标转写为视觉变体或可逆编码，分词切分通过打散词元边界破坏词法识别，低资源语言则将请求置于对齐覆盖相对薄弱的语言分布之下。这类攻击主要用于检验安全机制对浅层文本特征的依赖程度\cite{zhao2024accelerating,jiang2024artprompt}\cite{deng2024multilingual,huang2026obscure}\cite{wallace2019triggers}。另一类作用于任务的语义框架，核心是改变模型对当前任务是什么的解读路径：格式伪装将危险请求包装为结构化数据对象，认知误导通过角色设定与场景推演重新定义任务性质，指令重构利用问答链与任务拆分掩盖危险目标的完整形态，上下文操控则借助 many-shot 前缀在单轮输入内构建影响判断的连贯语境\cite{zeng2024johnny,chao2024pair}\cite{anil2024manyshot,li2024drattack}。

\begin{table}[H]
\centering
\footnotesize
\setlength{\tabcolsep}{4pt}
\renewcommand{\arraystretch}{0.95}
\caption{静态测试集九个攻击分组构成统计表}
\label{tab:static-groups}
\begin{tabular}{p{0.22\textwidth}p{0.12\textwidth}p{0.12\textwidth}p{0.40\textwidth}}
\toprule
\textbf{攻击维度} & \textbf{样本数} & \textbf{占比} & \textbf{代表构造方式} \\
\midrule
对抗后缀 & 720 & 11.24\% & 标点扰动、Emoji 后缀、Unicode 变体 \\
ASCII 伪装 & 720 & 11.24\% & ASCII Art、字符图案、外观重写 \\
编码扰动 & 720 & 11.24\% & Base64、ROT13、替换编码 \\
低资源语言 & 646 & 10.08\% & 低资源语言翻译、跨语种表达变体 \\
格式伪装 & 720 & 11.24\% & JSON、XML、Markdown \\
分词切分 & 720 & 11.24\% & 插入分隔符、局部切分 \\
认知误导 & 720 & 11.24\% & 角色设定、危机推演、机制分析 \\
指令重构 & 720 & 11.24\% & 问答链、任务拆分、最后一问注入 \\
上下文操控 & 720 & 11.24\% & many-shot 前缀、长上下文、背景嵌入 \\
\bottomrule
\end{tabular}
\end{table}

\section{动态攻击扩展：自动化红队与多轮越狱机制}

静态测试集有两个难以回避的局限：固定模板无法模拟真实攻击者反复改写提示的行为，单轮评估也看不到安全边界在持续对话中被逐步侵蚀的过程。为此，本文在静态框架之外另行构建了自动化红队与多轮越狱两个扩展模块，分别对应这两个方向上的不足。

\subsection{自动化红队攻击变体生成}

真实攻击者遭遇拒答后往往不会就此放弃，而是持续改写表达，直至触碰安全约束的薄弱点。自动化红队模块以此为出发点，借助辅助语言模型模拟这一改写行为：给定危险种子提示，自动生成若干语义近似但措辞各异的攻击变体，扩充单轮测试的覆盖空间\cite{perez2022redteaming,chao2024pair}\cite{liu2023autodan,mehrotra2024tap}。

流程上，模块依次完成种子加载、模板组装、代理模型变体生成、质量门筛选与结果入库。在变体生成之前，参考选例阶段先从全局样本池中为每条种子提示挑选 3 至 5 条风格锚点，选取标准融合攻击模式与词汇相似两个维度：
\begin{equation}
s_{\text{ref}}(x, r) = 0.4 \cdot \text{PatternOverlap}(x, r) + 0.6 \cdot \text{LexicalOverlap}(x, r),
\end{equation}
$\text{PatternOverlap}$ 从角色扮演、开发者模式、指令忽视等攻击策略维度衡量目标与候选的共性，$\text{LexicalOverlap}$ 则在词项层面补充语义相似度。二者加权组合，使选出的参考样本在攻击类型与表达风格上都尽量贴近种子提示。

生成变体须通过加权质量门方可写入结果：
\begin{equation}
\mathcal{Q}(x', x) = 0.55 \cdot S_{\text{intent}}(x', x) + 0.40 \cdot S_{\text{anchor}}(x', x) + 0.05 \cdot S_{\text{template}}(x', x),
\end{equation}
其中 $S_{\text{intent}}$ 基于中英文 bigram 的 Jaccard 相似度衡量语义保真度，$S_{\text{anchor}}$ 统计原始提示中意图锚词（中文 2--3 字词组，英文 4 字以上词项）在变体中的命中率，$S_{\text{template}}$ 考察风格符合程度。三项得分任一未达阈值（综合 0.45、意图 0.35、锚词覆盖 0.50）时，模块最多发起三轮重试，仍不达标则标记失败。通过质量门的变体以 JSONL 格式写出，记录字段与主实验管线保持一致，可直接加载执行。

\subsection{多轮交互式越狱}

多轮越狱模块将越狱过程建模为跨多轮的动态博弈：攻击者在每轮回复后持续调整策略，观察模型安全边界是否随对话历史的积累而松动\cite{anil2024manyshot,chao2024pair}。这与第二章讨论的上下文依赖漏洞相呼应，即便单条提示上的安全识别相对稳定，连续对话构建的语境本身也可能影响模型对同一危险目标的判定。

在实现上，模块首先由规划模型为整个 $K$ 轮交互预生成一份结构化攻击计划，约定每轮的攻击目标、策略方向与具体提示候选，以此规避逐轮临时生成容易引发的目标漂移问题。每轮交互结束后，判定模型输出结构化反馈，指出当前拒答所处阶段、可行的策略调整方向以及应规避的重复路径；下一轮提示据此生成，确保策略演进具有明确依据而非凭借随机试探。当连续两轮均遭拒答时，系统触发重规划流程，将累积的对话历史与判定反馈一并提交规划模型，重新生成剩余轮次的攻击计划。

每轮循环的逻辑较为固定：初始提示发出后，在线判定器即时评估回复——判定成功则终止，否则依据结构化反馈调整策略并生成追问，进入下一轮。每轮结果实时写入日志，字段涵盖当轮提示、模型回复、判定状态及当前策略。此外，模块在输入、交互与输出三个层面均预留了防御接口，同一套框架可同时用于攻击侧与防御侧的评估。输出指标除常规越狱成功率外，还包含首次越狱成功轮次与重规划触发频率，这类时序指标能够反映安全边界随对话积累的动态变化。

\section{多模态越狱攻击扩展讨论}

本文框架以文本输入为主要作用域，但视觉—语言模型（Vision-Language Models，VLMs）的普及使这一边界逐渐受到挑战。VLMs 将图像编码器与语言模型对接，而现有安全对齐机制通常只覆盖语言侧，视觉通道本身缺乏等效的约束，攻击者因此获得了绕过纯文本安全检测的新路径\cite{qi2024visualadversarial,shayegani2024jailbreakpieces}。

从机理角度看，多模态越狱大致可对应到本文九类文本攻击维度中的三个方向。

\textbf{对抗性图像注入。}这类攻击在图像上叠加人眼难以察觉的扰动，诱使 VLM 在视觉编码阶段将图像内容误读为特定危险语义，从而在文本输出侧产生违规内容\cite{qi2024visualadversarial}。其思路与编码扰动、ASCII 伪装等表征层攻击相近——同样是改变输入的表征形式以规避安全检测，只是扰动的载体从文本转移到了图像。

\textbf{排版视觉越狱。}攻击者将危险指令以排版文字的形式印入图片（如截图或图表），模型在视觉理解阶段直接读取并响应这些指令，而语言侧的关键词过滤因内容未经文本通道传入而完全失效。FigStep\cite{gong2023figstep} 是该方向的代表性工作。这与格式伪装、分词切分的逻辑相通，本质上都是借助载体形式的转换来隐藏危险内容。

\textbf{跨模态语义包装。}攻击者将危险请求拆分到图像和文本两个模态中：图像负责提供危险前提或语境，文本只呈现一个看似无害的问题，两部分单独审查均不触发拦截，合并后才构成完整的攻击意图\cite{shayegani2024jailbreakpieces}。这与认知误导、上下文操控的机理一脉相承，分步拆解、合并触发的思路从单模态延伸到了跨模态场景。

从以上分析可以看出，多模态越狱并未带来本质上全新的攻击原理，而是将文本层面的既有机理迁移到了视觉通道，本文框架的维度划分在结构上与这些变体相容。相关评测基准已有初步工作\cite{liu2024mmsafety}，其核心评估逻辑在配套 VLM 数据集的条件下可以复用。

\section{本章小结}

本章以黑盒威胁模型为基本假设，完成了多维越狱攻击框架从问题定义到具体构建的整体设计。

本章首先给出越狱攻击的操作性定义与量化指标，确立了贯穿全文实验的评估基准。测试数据覆盖多个风险类别，并纳入良性请求样本，兼顾对危险请求突破率与正常请求误判率的双向核验。

框架由静态测试集、自动化红队与多轮越狱三类攻击来源构成，共享同一执行管线以确保横向比较的口径一致。静态测试集规模较大，九个攻击分组按作用层次分为表征编码与语义框架两类：前者通过改变输入的呈现形式检验安全机制对浅层特征的依赖，后者介入任务意图的解读路径考察语义包装手段的干扰效果。自动化红队模块生成多样化攻击变体，模拟攻击者遭遇拒答后反复改写的行为；多轮越狱模块将越狱建模为动态博弈，量化安全边界在连续对话中的演化过程。

此外，本章就视觉—语言模型场景作了机理层面的扩展分析，说明框架的维度划分具备向多模态场景延伸的结构基础。以上内容共同构成第四章实验评测的方法依据。
